
\section{Conclusion}

We have systematically investigated whether Large Language Models can perform genuine analytical discovery. Our answer is nuanced: \textbf{LLMs excel at reproducing known patterns but fail catastrophically on specialized domains requiring extrapolation beyond training data.}

Theorem~\ref{thm:llm_reproductive} provides theoretical grounding: autoregressive generation probability decays exponentially for expressions distant from training corpora. Empirical validation across 131 tests confirms this, showing 95\% success on classical formulas versus 45\% on decentralized finance, with 412-847\% extrapolation errors.

\textsc{HypatiaX} addresses these limitations through careful role allocation: symbolic methods perform mathematical reasoning under physics constraints, while LLMs provide natural language interpretation and optional warm-starting. This hybrid approach achieves 88.9-100\% success rates with 8-23\% extrapolation error and 100\% functional correctness on ground truth tests.

Classical learning theory \citep{vapnik1999overview,bishop2006pattern,goodfellow2016deep} provides foundations for understanding generalization, though modern overparameterized models challenge traditional assumptions. Recent comprehensive surveys \citep{zhang2025survey} of reasoning LLMs confirm the need for hybrid approaches. The future of analytical discovery is not "AI vs. humans" or "neural vs. symbolic" but rather \textbf{intelligent orchestration of complementary capabilities}: LLMs as interfaces to symbolic discovery engines, with rigorous validation ensuring reliability.

\section*{Acknowledgments}

The author thanks domain experts who participated in blind evaluation of generated reports, and the open-source communities behind PySR, SymPy, and Anthropic Claude.

\bibliography{bibliography}

\appendix

\section{Complete Proof of Theorem~\ref{thm:llm_reproductive}}
\label{app:proof}

\begin{proof}[Detailed Proof]
Let $M_\theta$ be an autoregressive language model with parameters $\theta$ trained on corpus $\mathcal{C}$. Expression $E = t_1 \cdots t_n$ is generated via:
\[
P_\theta(E \mid D) = \prod_{i=1}^n P_\theta(t_i \mid t_{<i}, D)
\]

where $t_{<i} = t_1, \ldots, t_{i-1}$ is the prefix and $D$ represents conditioning data.

\paragraph{Step 1: Training concentrates probability mass.}

Maximum likelihood training optimizes:
\[
\theta^* = \arg\max_\theta \mathbb{E}_{E \sim \mathcal{C}}[\log P_\theta(E)]
\]

This causes $P_{\theta^*}(t_i \mid t_{<i}, D)$ to concentrate on tokens frequently observed in $\mathcal{C}$ given context $t_{<i}$.

\paragraph{Step 2: Novel tokens incur exponential penalty.}

For expression $E \notin \mathcal{C}$, define:
\[
d(E, \mathcal{C}) = \min_{E' \in \mathcal{C}} d_{\text{edit}}(E, E')
\]

Let $\mathcal{N}(E) \subseteq \{1, \ldots, n\}$ denote positions where $t_i$ is novel (low probability given $t_{<i}$). Under independence approximation:
\[
P_\theta(E \mid D) \approx \prod_{i \in \mathcal{N}(E)} p_i \cdot \prod_{i \notin \mathcal{N}(E)} p_i'
\]

where $p_i \ll 1$ for novel tokens and $p_i' \approx 1$ for familiar tokens.

\paragraph{Step 3: Exponential decay.}

Since $|\mathcal{N}(E)| \geq d(E, \mathcal{C})$ and $p_i \leq p_{\max} < 1$:
\[
P_\theta(E \mid D) \leq p_{\max}^{d(E, \mathcal{C})} = \exp(-\beta \cdot d(E, \mathcal{C}))
\]

where $\beta = -\log p_{\max} > 0$. Thus generation probability decays exponentially with edit distance from training data, establishing reproductive behavior. \qed
\end{proof}

\paragraph{Remark on rigor:} This proof assumes approximate local independence of token probabilities. A fully rigorous treatment would analyze the joint distribution over all tokens, but empirical evidence strongly supports the conclusion even under this approximation.

\section{Hyperparameter Settings}
\label{app:hyperparameters}

\begin{table}[H]
\centering
\caption{PySR Configuration}
\label{tab:hyperparameters}
\begin{tabular}{@{}lc@{}}
\toprule
\textbf{Parameter} & \textbf{Value} \\
\midrule
Population size & 100 \\
Generations & 500 \\
Populations (parallel) & 12 \\
Crossover probability & 0.8 \\
Mutation probability & 0.1 \\
Complexity penalty & 0.01 \\
Tournament size & 3 \\
Elitism & 5 \\
Max expression depth & 15 \\
Optimization iterations & 50 \\
\midrule
\textbf{Validation Thresholds} & \\
\midrule
Acceptance $R^2$ (Criterion 1) & 0.99 \\
Validation score (Criterion 1) & 30 \\
Acceptance $R^2$ (Criterion 2) & 0.95 \\
Validation score (Criterion 2) & 80 \\
Extrapolation $R^2$ & 0.80 \\
Numerical stability ($\alpha$) & 0.95 \\
\bottomrule
\end{tabular}
\end{table}

\section{Detailed Experimental Results}
\label{app:detailed_results}

\subsection{Pure LLM Baseline (40 tests)}

\begin{table}[H]
\centering
\caption{Pure LLM Performance by Test Case}
\label{tab:llm_detailed}
\begin{small}
\begin{tabular}{@{}llccp{4.5cm}@{}}
\toprule
\textbf{Domain} & \textbf{Test} & \textbf{$R^2$} & \textbf{Result} & \textbf{Failure Mode} \\
\midrule
\multicolumn{5}{c}{\textit{Classical Scientific Domains (20 tests)}} \\
\midrule
Materials & Stress-strain & 1.00 & Pass & --- \\
Materials & Young's modulus & 1.00 & Pass & --- \\
Materials & Poisson ratio & 1.00 & Pass & --- \\
Materials & Thermal expansion & 1.00 & Pass & --- \\
Fluid & Reynolds number & 1.00 & Pass & --- \\
Fluid & Bernoulli equation & 1.00 & Pass & --- \\
Fluid & Drag coefficient & 1.00 & Pass & --- \\
Fluid & Flow rate & 1.00 & Pass & --- \\
Thermo & Ideal gas law & 1.00 & Pass & --- \\
Thermo & Heat capacity & 1.00 & Pass & --- \\
Thermo & Carnot efficiency & 1.00 & Pass & --- \\
Thermo & Stefan-Boltzmann & 1.00 & Pass & --- \\
Mechanics & Newton's 2nd law & 1.00 & Pass & --- \\
Mechanics & Kinetic energy & 1.00 & Pass & --- \\
Mechanics & Gravitational force & 1.00 & Pass & --- \\
Mechanics & Hooke's law & 1.00 & Pass & --- \\
Chemistry & Arrhenius equation \citep{arrhenius1889reaction} & 1.00 & Pass & --- \\
Chemistry & pH calculation & 1.00 & Pass & --- \\
Chemistry & Nernst equation & 0.50 & Fail & Wrong sign convention \\
Chemistry & Rate law & 1.00 & Pass & --- \\
\midrule
\multicolumn{5}{c}{\textit{Decentralized Finance Domains (20 tests)}} \\
\midrule
AMM & Constant product & 1.00 & Pass & --- \\
AMM & Price impact & 0.98 & Pass & --- \\
AMM & Impermanent loss & $-1.26$ & Fail & Unit inconsistency (missing $\times$100) \\
AMM & Swap fee & 0.82 & Pass & --- \\
VaR & Parametric VaR 95\% & $-10.05$ & Fail & Wrong sign (z-score) \\
VaR & Parametric VaR 99\% & 1.00 & Pass & --- \\
VaR & Historical VaR & 0.95 & Pass & --- \\
VaR & Modified VaR & $-2.15$ & Fail & Scipy API misuse \\
Liquidity & APY calculation & 1.00 & Pass & --- \\
Liquidity & Capital efficiency & 0.87 & Pass & --- \\
Liquidity & Fee earnings & --- & Fail & Non-executable (tutorial mode) \\
Liquidity & Utilization ratio & --- & Fail & Non-executable (tutorial mode) \\
ES & Expected Shortfall 95\% & $-4.12$ & Fail & Scipy \texttt{norm.pdf} returns object \\
ES & Expected Shortfall 99\% & $-3.87$ & Fail & Distribution tail error \\
ES & Conditional VaR & 0.88 & Pass & --- \\
ES & Tail risk & $-1.92$ & Fail & Statistical definition error \\
Liquidation & Long position & --- & Fail & Tutorial text only \\
Liquidation & Short position & --- & Fail & Tutorial text only \\
Liquidation & Margin call & --- & Fail & Incomplete derivation \\
Liquidation & Leverage ratio & --- & Fail & Multi-step missing \\
\midrule
\multicolumn{3}{l}{\textbf{Classical Average:}} & \textbf{19/20 (95\%)} & \\
\multicolumn{3}{l}{\textbf{DeFi Average:}} & \textbf{8/20 (40\%)} & \\
\multicolumn{3}{l}{\textbf{Overall:}} & \textbf{27/40 (67.5\%)} & \\
\bottomrule
\end{tabular}
\end{small}
\end{table}

\subsection{DeFi-Optimized Results (23 tests)}

\begin{table}[H]
\centering
\caption{DeFi-Optimized HypatiaX Performance}
\label{tab:defi_detailed}
\begin{small}
\begin{tabular}{@{}lccccc@{}}
\toprule
\textbf{Test} & \textbf{$R^2$} & \textbf{Validation} & \textbf{Time (s)} & \textbf{Extrap.} & \textbf{Result} \\
\midrule
Constant product (xy=k) & 1.0000 & 100.0 & 45 & Pass & Pass \\
Price impact & 0.9998 & 95.2 & 120 & Pass & Pass \\
Impermanent loss & 0.9995 & 92.1 & 180 & Pass & Pass \\
Swap fee calculation & 1.0000 & 98.5 & 60 & Pass & Pass \\
VaR 95\% parametric & 0.9988 & 88.3 & 150 & Pass & Pass \\
VaR 99\% parametric & 0.9990 & 90.1 & 145 & Pass & Pass \\
Historical VaR & 0.9985 & 87.5 & 135 & Pass & Pass \\
Modified VaR & 0.9982 & 85.9 & 160 & Pass & Pass \\
APY simple & 1.0000 & 100.0 & 50 & Pass & Pass \\
APY compound & 0.9995 & 94.8 & 110 & Pass & Pass \\
Capital efficiency & 0.9940 & 78.5 & 240 & Pass & Pass \\
Fee earnings & 0.9998 & 96.2 & 90 & Pass & Pass \\
Expected Shortfall 95\% & 0.9987 & 89.7 & 170 & Pass & Pass \\
Expected Shortfall 99\% & 0.9985 & 88.2 & 175 & Pass & Pass \\
Conditional VaR & 0.9990 & 91.4 & 155 & Pass & Pass \\
Tail risk metric & 0.9983 & 86.8 & 165 & Pass & Pass \\
Liquidation (long) & 0.9992 & 93.6 & 125 & --- & Pass \\
Liquidation (short) & 0.9994 & 94.1 & 120 & --- & Pass \\
Margin call threshold & 0.9996 & 95.8 & 105 & --- & Pass \\
Effective leverage & 0.9998 & 97.2 & 85 & --- & Pass \\
Staking rewards & 1.0000 & 100.0 & 55 & --- & Pass \\
Dilution effect & 0.9993 & 93.8 & 115 & --- & Pass \\
Kelly criterion & 0.9988 & 89.5 & 195 & Pass & Pass \\
\midrule
\textbf{Average} & \textbf{0.9990} & \textbf{91.3} & \textbf{115.2} & \textbf{6/6} & \textbf{23/23} \\
\bottomrule
\end{tabular}
\end{small}
\end{table}

\paragraph{Extrapolation Tests:} Six tests designated for extrapolation validation (constant product, impermanent loss, VaR variants, ES variants, Kelly criterion) all passed with average extrapolation $R^2 = 0.9650$, demonstrating robust generalization.

\section{Reproducibility Statement}
\label{app:reproducibility}

All experimental data, source code, and configuration files are publicly available to ensure full reproducibility of our results.

\subsection{Code Repository}

\textbf{GitHub:} \url{https://github.com/ruperto-bonet/hybrid-symbolic-discovery}

\paragraph{Repository Structure:}
\begin{small}
\begin{verbatim}
hybrid-symbolic-discovery/
|-- hypatiax/
|   |-- core/
|   |   |-- generation/         # Discovery engines
|   |   |-- validation/         # Multi-layer validators
|   |-- tools/
|   |   |-- symbolic/           # PySR integration
|   |-- data/
|       |-- results/            # Session outputs
|-- experiments/
|   |-- pure_llm_baseline.py    # 40 LLM tests
|   |-- protocol_a.py           # 18 symbolic tests
|   |-- protocol_all_30.py      # LLM-guided hybrid
|   |-- defi_suite.py           # 23 DeFi tests
|-- notebooks/
|   |-- analysis.ipynb          # Result visualization
|-- docker/
    |-- Dockerfile              # Reproducible environment
\end{verbatim}
\end{small}

\subsection{Session Data}

All raw results are stored with timestamps:

\begin{itemize}
\item \textbf{Pure LLM Baseline}: \texttt{data/results/pure\_llm\_20260115/}
\item \textbf{Protocol A (18 tests)}: \texttt{data/results/20260114\_153042/}
\item \textbf{LLM-Guided (30 tests)}: \texttt{data/results/llm\_20260114\_183940/}
\item \textbf{DeFi Suite (23 tests)}: \texttt{data/results/defi\_20260113\_142837/}
\end{itemize}

Each session contains:
\begin{itemize}
\item Discovered expressions (LaTeX and Python)
\item Performance metrics (R², validation scores, timing)
\item Full validation reports (dimensional analysis, constraint checking)
\item Extrapolation test results
\end{itemize}

\subsection{Running Experiments}

\paragraph{Quick Start:}
\begin{small}
\begin{verbatim}
# Clone repository
git clone https://github.com/ruperto-bonet/hybrid-symbolic-discovery
cd hybrid-symbolic-discovery

# Install dependencies
pip install -r requirements.txt

# Run pure symbolic baseline (18 tests, ~2 hours)
python experiments/protocol_a.py --batch --mode STANDARD

# Run LLM-guided hybrid (30 tests, ~35 minutes)
python experiments/protocol_all_30.py --batch --mode hybrid

# Run DeFi suite (23 tests, ~45 minutes)
python experiments/defi_suite.py --batch --mode STANDARD
\end{verbatim}
\end{small}

\paragraph{Resuming Interrupted Runs:}
\begin{small}
\begin{verbatim}
python experiments/protocol_all_30.py --resume --session-id 20260114_183940
\end{verbatim}
\end{small}

\subsection{Hardware Requirements}

\begin{itemize}
\item \textbf{CPU}: 8+ cores recommended (tested on AMD Ryzen 9 5950X)
\item \textbf{RAM}: 16GB minimum, 32GB+ recommended
\item \textbf{Storage}: 10GB for code + results
\item \textbf{OS}: Linux (Ubuntu 22.04), macOS, or Windows with WSL2
\item \textbf{Python}: 3.9+
\end{itemize}

\subsection{Software Dependencies}

\begin{small}
\begin{verbatim}
PySR==0.16.3           # Symbolic regression
SymPy==1.12            # Symbolic mathematics
NumPy==1.24.3          # Numerical computing
Pandas==2.0.3          # Data handling
SciPy==1.11.2          # Scientific computing
Matplotlib==3.7.2      # Visualization
Anthropic==0.18.1      # LLM API (optional for LLM-guided mode)
\end{verbatim}
\end{small}

\subsection{Random Seeds}

All experiments use fixed random seeds for reproducibility:

\begin{itemize}
\item Pure symbolic: \texttt{seed=42}
\item LLM-guided: \texttt{seed=42} (PySR component)
\item DeFi suite: \texttt{seed=42}
\end{itemize}

LLM outputs have inherent variability due to temperature sampling (temperature=0.7), but validation ensures only mathematically correct expressions are accepted regardless of generation randomness.

\subsection{Expected Runtime}

On reference hardware (Ryzen 9 5950X, 32GB RAM):

\begin{table}[H]
\centering
\begin{tabular}{@{}lcc@{}}
\toprule
\textbf{Experiment} & \textbf{Tests} & \textbf{Wall Time} \\
\midrule
Pure LLM Baseline & 40 & 5-10 min \\
Protocol A (Pure Symbolic) & 18 & 117 min \\
LLM-Guided Hybrid & 30 & 35 min \\
DeFi Suite & 23 & 44 min \\
Extended Validation & 20 & 136 min \\
\midrule
\textbf{All Experiments} & \textbf{131} & \textbf{~5.5 hours} \\
\bottomrule
\end{tabular}
\end{table}

\subsection{Verification Checksums}

We provide SHA-256 checksums for critical result files:

\begin{small}
\begin{verbatim}
# Verify result integrity
sha256sum data/results/protocol_a_summary.json
# Expected: 8f3a2b1c... (provided in repository)
\end{verbatim}
\end{small}

\subsection{Docker Container}

For maximum reproducibility, we provide a Docker container with all dependencies pre-configured:

\begin{small}
\begin{verbatim}
# Build container
docker build -t hypatiax:v1.0 -f docker/Dockerfile .

# Run experiments
docker run -v $(pwd)/results:/results hypatiax:v1.0 \
    python experiments/protocol_all_30.py --batch
\end{verbatim}
\end{small}

\subsection{Contact for Support}

For issues with reproduction:
\begin{itemize}
\item Open GitHub issue: \url{https://github.com/ruperto-bonet/hybrid-symbolic-discovery/issues}
\item Email: \texttt{ruperto.bonet@modelphysmat.com}
\end{itemize}

We commit to maintaining the repository and responding to reproduction issues within 2 weeks.

\section{Computational Complexity Analysis}
\label{app:complexity}

\subsection{Theoretical Complexity}

\begin{proposition}[HypatiaX Complexity]
Let $n$ be candidate expressions, $m$ evaluation points, $k$ symbolic constraints. Total complexity is $\mathcal{O}(n \cdot (k + m))$.
\end{proposition}

\begin{proof}
For each candidate expression $E$:
\begin{enumerate}
\item \textbf{Symbolic constraint checking}: $\mathcal{O}(k)$ operations (dimensional analysis, domain rules)
\item \textbf{Numerical evaluation}: $\mathcal{O}(m)$ operations (fitness computation over dataset)
\end{enumerate}

Total per candidate: $\mathcal{O}(k + m)$. With $n$ candidates: $\mathcal{O}(n(k + m))$. \qed
\end{proof}

\subsection{Empirical Timing Breakdown}

We profiled HypatiaX on representative test cases:

\begin{table}[H]
\centering
\caption{Computational Cost Breakdown (avg per test)}
\label{tab:timing_breakdown}
\begin{tabular}{@{}lcc@{}}
\toprule
\textbf{Component} & \textbf{Time (s)} & \textbf{Percentage} \\
\midrule
\multicolumn{3}{c}{\textit{Pure Symbolic (PySR)}} \\
\midrule
Population initialization & 5.2 & 1.3\% \\
Evolutionary search & 365.8 & 93.9\% \\
Symbolic simplification & 12.3 & 3.2\% \\
Validation framework & 6.7 & 1.6\% \\
\midrule
\textbf{Total} & \textbf{390.0} & \textbf{100\%} \\
\midrule
\multicolumn{3}{c}{\textit{LLM-Guided Hybrid}} \\
\midrule
LLM API call & 8.5 & 12.1\% \\
Expression parsing & 2.0 & 2.9\% \\
Initial validation & 3.0 & 4.3\% \\
PySR refinement (when needed) & 52.5 & 75.0\% \\
Final validation & 4.0 & 5.7\% \\
\midrule
\textbf{Total} & \textbf{70.0} & \textbf{100\%} \\
\bottomrule
\end{tabular}
\end{table}

\paragraph{Key Findings:}
\begin{itemize}
\item Evolutionary search dominates (93.9\%) in pure symbolic mode
\item LLM-guided mode reduces this by solving 73\% of cases without PySR
\item Validation overhead is negligible (1.6-5.7\%)
\end{itemize}

\subsection{Scalability Analysis}

We tested scalability with varying dataset sizes:

\begin{table}[H]
\centering
\caption{Scalability with Dataset Size}
\label{tab:scalability}
\begin{tabular}{@{}cccc@{}}
\toprule
\textbf{Data Points} & \textbf{Pure PySR (s)} & \textbf{LLM-Hybrid (s)} & \textbf{Speedup} \\
\midrule
100 & 245 & 48 & 5.1$\times$ \\
500 & 390 & 70 & 5.6$\times$ \\
1000 & 521 & 95 & 5.5$\times$ \\
5000 & 1245 & 215 & 5.8$\times$ \\
\bottomrule
\end{tabular}
\end{table}

Speedup remains relatively constant across dataset sizes, confirming that LLM warm-starting provides consistent efficiency gains.


\section*{Supplementary Materials}

Comprehensive empirical analysis of LLM performance across 140 formula generations and 10 scientific domains is available in the supplementary materials document. This includes:

\begin{itemize}
\item Detailed failure mode taxonomy with code examples
\item Domain-specific performance analysis (Materials Science, Fluid Dynamics, Thermodynamics, Mechanics, Chemistry, DeFi)
\item Production deployment framework with validation pipelines
\item Complete experimental logs from 7 runs (Run IDs: 162102, 163240, 165852, 173905, 192807, 232733, 235258)
\end{itemize}

Key supplementary findings:
\begin{itemize}
\item Pure LLM success variance of 0-100\% across runs using identical test cases
\item Perfect symbolic accuracy ($R^2 = 1.0000$) achieved on 75\% of formulas when prompts are optimized
\item Complete system failures (0\% success rate) in 3/7 runs due to API errors, variable binding issues, and code extraction failures
\item Prompt engineering quality is the dominant factor determining success, capable of transforming catastrophic failures into production-ready implementations
\end{itemize}

\section{Extended Discussion}
\label{app:extended_discussion}

\subsection{Why LLMs Fail on Novel Domains}

Our empirical results (95\% classical vs. 45\% DeFi) raise a question: \textbf{Why do LLMs fail specifically on specialized domains?}

\paragraph{Training Data Coverage:}

Classical physics formulas appear extensively in:
\begin{itemize}
\item Textbooks (digitized and crawled)
\item Wikipedia and educational websites
\item Academic papers (often with equation reproduction in intro sections)
\item Code repositories (physics simulations, homework solutions)
\item Q\&A forums (StackExchange, Quora)
\end{itemize}

DeFi formulas, by contrast:
\begin{itemize}
\item Emerged post-2020 (after many LLM training cutoffs)
\item Appear primarily in white papers and technical documentation
\item Often use non-standard notation
\item Require understanding of cryptographic primitives and game theory
\item Limited presence in traditional academic literature
\end{itemize}

\paragraph{Conceptual Complexity:}

\begin{table}[H]
\centering
\caption{Conceptual Requirements Comparison}
\label{tab:conceptual_complexity}
\begin{tabular}{@{}lcc@{}}
\toprule
\textbf{Requirement} & \textbf{Classical Physics} & \textbf{DeFi/Risk} \\
\midrule
Distributional reasoning & Rare & Essential \\
Multi-step derivations & Common but formulaic & Domain-specific \\
Statistical API knowledge & Basic & Advanced (scipy.stats) \\
Convention awareness & Well-established & Evolving \\
Asymptotic analysis & Standard & Context-dependent \\
\bottomrule
\end{tabular}
\end{table}

\subsection{Lessons for Future LLM Development}

Our findings suggest design principles for more reliable scientific LLMs:

\begin{enumerate}
\item \textbf{Training data curation}: Include specialized domain corpora with verified formulas and extensive code examples

\item \textbf{Execution-based feedback}: Train with reward signals from actual code execution, not just text likelihood

\item \textbf{Explicit uncertainty quantification}: Models should output confidence scores calibrated to actual correctness probability

\item \textbf{Symbolic grounding}: Integrate symbolic math engines (SymPy, Mathematica) during inference for dimensional checking

\item \textbf{Retrieval augmentation}: Query domain-specific databases (NIST, financial standards) before generation
\end{enumerate}

Even with these improvements, Theorem~\ref{thm:llm_reproductive} suggests fundamental limits on extrapolation capability.

\subsection{Broader Applicability}

While we focused on physics and finance, the hybrid architecture generalizes to other domains requiring analytical discovery:

\paragraph{Candidate Domains:}
\begin{itemize}
\item \textbf{Materials science}: Stress-strain relationships, phase diagrams
\item \textbf{Pharmacokinetics}: Drug absorption, distribution, metabolism models
\item \textbf{Ecology}: Population dynamics, predator-prey systems
\item \textbf{Economics}: Production functions, utility optimization
\item \textbf{Climate science}: Radiative forcing, feedback mechanisms
\end{itemize}

\paragraph{Required Adaptations:}
\begin{itemize}
\item Domain-specific operator sets
\item Discipline-appropriate constraint libraries
\item Validation metrics aligned with domain standards
\item Expert-in-the-loop refinement for unfamiliar domains
\end{itemize}

\subsection{Limitations and Future Directions}

\paragraph{Current Limitations:}

\begin{enumerate}
\item \textbf{Operator completeness}: If true formula requires operators outside $\mathcal{O}$, discovery fails
\item \textbf{Multi-scale phenomena}: Expressions with vastly different characteristic scales challenge numerical evaluation
\item \textbf{Implicit constraints}: Some domain knowledge is tacit and difficult to formalize
\item \textbf{Computational cost}: Evolutionary search remains expensive for complex problems
\end{enumerate}

\paragraph{Promising Directions:}

\begin{enumerate}
\item \textbf{Active learning}: Use LLM feedback to guide symbolic search more intelligently:
\begin{verbatim}
While not converged:
    E* = PySR.search()
    suggestions = LLM.critique(E*)
    constraints.add(suggestions)
\end{verbatim}

\item \textbf{Hierarchical discovery}: Discover sub-expressions separately, then combine
\item \textbf{Transfer learning}: Leverage formulas from related domains
\item \textbf{Uncertainty quantification}: Provide confidence intervals on discovered expressions
\item \textbf{Multi-modal integration}: Incorporate experimental images, diagrams, protocol descriptions
\end{enumerate}

\end{document}
